%% ============================================================
%%  Stock Tracking System — Technical Paper
%%  Author  : Alper
%%  Date    : April 2026
%%  Engine  : pdfLaTeX
%% ============================================================
\documentclass[12pt,a4paper,twoside]{article}

%% ---- Encoding & Language ----
\usepackage[utf8]{inputenc}
\usepackage[T1]{fontenc}
\usepackage[turkish,english]{babel}   % English primary, Turkish support

%% ---- Typography ----
\usepackage{lmodern}
\usepackage{microtype}
\usepackage{setspace}
\onehalfspacing

%% ---- Page Geometry ----
\usepackage[
  a4paper,
  top=2.5cm, bottom=2.5cm,
  left=3cm,  right=2.5cm,
  headheight=15pt
]{geometry}

%% ---- Headers & Footers ----
\usepackage{fancyhdr}
\pagestyle{fancy}
\fancyhf{}
\fancyhead[LE,RO]{\small\thepage}
\fancyhead[RE]{\small\textit{Stock Tracking System}}
\fancyhead[LO]{\small\textit{\leftmark}}
\renewcommand{\headrulewidth}{0.4pt}

%% ---- Graphics & Color ----
\usepackage{graphicx}
\usepackage{xcolor}
\usepackage{tikz}
\usepackage{pgf-pie}
\usetikzlibrary{
  shapes.geometric, shapes.multipart, shapes.symbols,
  arrows.meta, positioning, fit, backgrounds,
  decorations.pathmorphing, calc, patterns, shadows
}

%% ---- Tables ----
\usepackage{booktabs}
\usepackage{tabularx}
\usepackage{multirow}
\usepackage{array}
\usepackage{longtable}

%% ---- Listings (code blocks) ----
\usepackage{listings}
\usepackage{lstautogobble}

\definecolor{codebg}{HTML}{F5F5F5}
\definecolor{codeframe}{HTML}{CCCCCC}
\definecolor{keyword}{HTML}{0033B3}
\definecolor{string}{HTML}{067D17}
\definecolor{comment}{HTML}{8C8C8C}
\definecolor{number}{HTML}{1750EB}

\lstdefinestyle{javaStyle}{
  backgroundcolor=\color{codebg},
  frame=single, rulecolor=\color{codeframe},
  basicstyle=\ttfamily\small,
  keywordstyle=\color{keyword}\bfseries,
  stringstyle=\color{string},
  commentstyle=\color{comment}\itshape,
  numberstyle=\tiny\color{number},
  numbers=left, numbersep=6pt,
  breaklines=true, autogobble=true,
  tabsize=2,
  language=Java
}

\lstdefinestyle{yamlStyle}{
  backgroundcolor=\color{codebg},
  frame=single, rulecolor=\color{codeframe},
  basicstyle=\ttfamily\small,
  keywordstyle=\color{keyword}\bfseries,
  stringstyle=\color{string},
  commentstyle=\color{comment}\itshape,
  breaklines=true, autogobble=true,
  tabsize=2,
  language=bash
}

\lstdefinestyle{sqlStyle}{
  backgroundcolor=\color{codebg},
  frame=single, rulecolor=\color{codeframe},
  basicstyle=\ttfamily\small,
  keywordstyle=\color{keyword}\bfseries,
  stringstyle=\color{string},
  commentstyle=\color{comment}\itshape,
  breaklines=true, autogobble=true,
  tabsize=2,
  language=SQL
}

%% ---- Math ----
\usepackage{amsmath, amssymb}

%% ---- Hyperlinks (last!) ----
\usepackage[
  colorlinks=true,
  linkcolor=black,
  citecolor=blue!70!black,
  urlcolor=blue!70!black,
  pdftitle={Stock Tracking System: A Distributed Architecture for CAP Theorem Demonstration},
  pdfauthor={Alper}
]{hyperref}

%% ---- Caption Formatting ----
\usepackage{caption}
\captionsetup{
  font=small, labelfont=bf,
  margin=0.5cm
}

%% ---- Miscellaneous ----
\usepackage{enumitem}
\usepackage{float}
\usepackage{placeins}
\usepackage{subcaption}
\usepackage{mdframed}

%% ============================================================
%%  Custom Colors & Macros
%% ============================================================
\definecolor{primary}{HTML}{1565C0}
\definecolor{secondary}{HTML}{0288D1}
\definecolor{accent}{HTML}{00ACC1}
\definecolor{success}{HTML}{2E7D32}
\definecolor{danger}{HTML}{C62828}
\definecolor{warning}{HTML}{F57F17}
\definecolor{neutral}{HTML}{546E7A}

\newcommand{\component}[1]{\textcolor{primary}{\textbf{#1}}}
\newcommand{\tech}[1]{\textit{#1}}
\newcommand{\cap}[1]{\textbf{\textcolor{primary}{#1}}}

%% ============================================================
%%  TikZ Node Styles
%% ============================================================
\tikzset{
  svcbox/.style={
    rectangle, rounded corners=4pt, draw=#1!80!black,
    fill=#1!20, text centered, minimum width=2.8cm,
    minimum height=0.8cm, font=\small\bfseries,
    drop shadow={shadow xshift=1pt, shadow yshift=-1pt,
                 fill=#1!50!black, opacity=0.3}
  },
  arrow/.style={-Stealth, thick, draw=neutral!70},
  bigbox/.style={
    rectangle, rounded corners=6pt,
    draw=neutral!60, fill=neutral!5,
    inner sep=8pt
  },
  layerlabel/.style={
    font=\bfseries\small, text=neutral
  },
  dbcyl/.style={
    cylinder, shape border rotate=90, draw=#1!80!black,
    fill=#1!20, aspect=0.3, minimum width=1.8cm,
    minimum height=1.2cm, font=\small\bfseries
  },
  capnode/.style={
    circle, draw=#1!80!black, fill=#1!20,
    minimum size=1.4cm, font=\small\bfseries,
    text centered, text width=1.2cm
  },
  msgbox/.style={
    rectangle, rounded corners=3pt,
    draw=primary!60, fill=primary!10,
    font=\ttfamily\small, inner sep=4pt
  }
}

%% ============================================================
\begin{document}

%% ---- Title Page ----
\begin{titlepage}
  \centering
  \vspace*{2cm}

  {\LARGE\bfseries\color{primary}
    Stock Tracking System\par}

  \vspace{0.6cm}
  {\Large\color{secondary}
    A Distributed Architecture for CAP Theorem Demonstration\par}

  \vspace{1.5cm}
  \rule{0.65\textwidth}{1.5pt}

  \vspace{1.5cm}
  {\large Alper\par}
  \vspace{0.3cm}
  {\normalsize April 2026\par}

  \vspace{3cm}

  %% Decorative architecture mini-diagram
  \begin{tikzpicture}[scale=0.85]
    \node[svcbox=secondary, minimum width=4cm] (fe)   at (0,0)   {React Frontend};
    \node[svcbox=primary,   minimum width=4cm] (be)   at (0,-1.8) {Spring Boot API};
    \node[svcbox=success,   minimum width=2cm] (db)   at (-3,-3.6){MySQL};
    \node[svcbox=accent,    minimum width=2cm] (rd)   at (0,-3.6) {Redis};
    \node[svcbox=warning,   minimum width=2cm] (kf)   at (3,-3.6) {Kafka};
    \node[svcbox=neutral,   minimum width=2cm] (pr)   at (-1.5,-5.2){Prometheus};
    \node[svcbox=neutral,   minimum width=2cm] (gr)   at (1.5,-5.2) {Grafana};

    \draw[arrow] (fe) -- (be);
    \draw[arrow] (be) -- (db);
    \draw[arrow] (be) -- (rd);
    \draw[arrow] (be) -- (kf);
    \draw[arrow] (be) -- (pr);
    \draw[arrow] (pr) -- (gr);
  \end{tikzpicture}

  \vfill
  {\small\color{neutral}
    Distributed Systems · Event Streaming · Fault Injection · IoT Integration}
\end{titlepage}

%% ---- Abstract ----
\begin{abstract}
\noindent
This paper presents the design, implementation, and evaluation of a
\textbf{distributed stock-tracking system} built to demonstrate the
\emph{CAP Theorem} in practice.
The system integrates a \tech{Spring Boot} backend, a \tech{Kafka}
event-streaming bus, a \tech{Redis} cache layer, a \tech{MySQL}
relational database, a \tech{Python} IoT sensor simulator, and a
\tech{React} single-page application — all orchestrated with
\tech{Docker Compose} across ten containerised services.
Network partitions are injected deterministically through
\tech{Toxiproxy}, allowing systematic observation of consistency/avail\-ability
trade-offs.
A \tech{Transactional Outbox Pattern} guarantees at-least-once event
delivery to Kafka during transient broker outages.
End-to-end observability is achieved via \tech{Prometheus} metrics and
dual \tech{Grafana} dashboards for stock operations and IoT sensor data.
Benchmarks demonstrate a stable throughput of 200 ops/s under normal
conditions, degrading predictably under injected latency, and maintaining
\emph{full functional availability} even when the cache or message broker
are artificially partitioned.
\end{abstract}

\tableofcontents
\newpage

%% ============================================================
\section{Introduction}
\label{sec:intro}
%% ============================================================

Modern enterprise applications are increasingly built on distributed
components that communicate over unreliable networks.
The \emph{CAP Theorem}~\cite{brewer2000towards} formalises an inherent
tension in such systems: it is impossible to guarantee
\textbf{Consistency (C)}, \textbf{Availability (A)}, and
\textbf{Partition Tolerance (P)} simultaneously.
Real systems must therefore make deliberate architectural choices —
favouring either \emph{CP} (consistency over availability) or
\emph{AP} (availability over consistency) during network partitions.

This project implements a production-like stock-tracking platform whose
primary educational and research objective is to \emph{make those
trade-offs observable and measurable}.
Each subsystem (\tech{MySQL}, \tech{Redis}, \tech{Kafka}) occupies a
different region of the CAP triangle; fault injection via
\tech{Toxiproxy} triggers real partition scenarios; and
\tech{Prometheus/Grafana} surfaces the resulting behaviour in real time.

\subsection{Objectives}
\begin{enumerate}[leftmargin=1.8cm, label=\textbf{O\arabic*.}]
  \item Design a representative distributed system covering data persistence,
        caching, event-streaming, and real-time IoT ingestion.
  \item Implement resilience patterns (Outbox, Cache-Aside, Rate Limiting)
        that trade consistency for availability where appropriate.
  \item Provide controlled, repeatable fault-injection experiments via
        Toxiproxy proxies.
  \item Surface system behaviour through 30+ Prometheus metrics and two
        Grafana dashboards (Stock Operations, IoT Sensor Metrics).
\end{enumerate}

\subsection{Paper Structure}

Section~\ref{sec:arch} describes the overall architecture.
Section~\ref{sec:components} details each sub-component.
Section~\ref{sec:dataflow} traces the principal data flows.
Section~\ref{sec:cap} analyses CAP positions and partition strategies.
Section~\ref{sec:observability} covers the monitoring stack.
Section~\ref{sec:iot} details the IoT pipeline.
Section~\ref{sec:testing} presents fault-injection and load-testing results.
Section~\ref{sec:design} discusses key design decisions and trade-offs.
Section~\ref{sec:conclusion} concludes with future work.

%% ============================================================
\section{System Architecture}
\label{sec:arch}
%% ============================================================

\subsection{High-Level Overview}

Figure~\ref{fig:hla} presents the three-layer architecture of the system.
The \textbf{User Layer} hosts the \tech{React 18} single-page application
served by an \tech{Nginx} reverse proxy.
The \textbf{Application Layer} contains the \tech{Spring Boot 3.x} backend,
which orchestrates all business logic.
The \textbf{Infrastructure Layer} provides persistent storage (\tech{MySQL 8}),
in-memory caching (\tech{Redis 7}), event streaming (\tech{Kafka 4.1 KRaft}),
metrics collection (\tech{Prometheus}), and visualisation (\tech{Grafana}).
\tech{Toxiproxy} sits transparently in front of each infrastructure
component, enabling deterministic network partition injection.

\begin{figure}[H]
\centering
\begin{tikzpicture}[node distance=0.6cm and 0.5cm, scale=0.88, every node/.style={scale=0.88}]

  %% ---- User Layer ----
  \begin{scope}[on background layer]
    \node[bigbox, fill=secondary!8, minimum width=12cm, minimum height=1.6cm]
      (userbox) at (0,0) {};
  \end{scope}
  \node[layerlabel, above left=0pt and 4pt of userbox.north west] {User Layer};
  \node[svcbox=secondary, minimum width=10cm] (fe) at (0,0)
    {React 18 + Vite \ (Nginx Reverse Proxy : 3000)};

  %% ---- Application Layer ----
  \begin{scope}[on background layer]
    \node[bigbox, fill=primary!8, minimum width=12cm, minimum height=1.6cm]
      (appbox) at (0,-2.2) {};
  \end{scope}
  \node[layerlabel, above left=0pt and 4pt of appbox.north west] {Application Layer};
  \node[svcbox=primary, minimum width=10cm] (be) at (0,-2.2)
    {Spring Boot 3.x \ (Java 17 · Gradle 8.5 : 8080)};

  %% ---- Infrastructure Layer ----
  \begin{scope}[on background layer]
    \node[bigbox, fill=neutral!6, minimum width=12cm, minimum height=2.4cm]
      (infbox) at (0,-5.0) {};
  \end{scope}
  \node[layerlabel, above left=0pt and 4pt of infbox.north west] {Infrastructure Layer};

  \node[svcbox=success, minimum width=2.2cm]  (mysql) at (-4.5,-4.8)  {MySQL\\:3306};
  \node[svcbox=accent,  minimum width=2.2cm]  (redis) at (-1.5,-4.8)  {Redis\\:6379};
  \node[svcbox=warning, minimum width=2.2cm]  (kafka) at (1.5,-4.8)   {Kafka\\:9092};
  \node[svcbox=neutral, minimum width=2.2cm]  (prom)  at (4.5,-4.8)   {Prometheus\\:9090};
  \node[svcbox=neutral, minimum width=2.2cm]  (graf)  at (4.5,-5.8)   {Grafana\\:3001};
  \node[svcbox=danger,  minimum width=2.2cm]  (toxi)  at (0,-5.8)     {Toxiproxy\\:8474};

  %% Arrows
  \draw[arrow] (fe)    -- node[right, font=\tiny]{Nginx /api/} (be);
  \draw[arrow] (be)    -- (mysql);
  \draw[arrow] (be)    -- (redis);
  \draw[arrow] (be)    -- (kafka);
  \draw[arrow] (be)    -- (prom);
  \draw[arrow] (prom)  -- (graf);
  \draw[arrow, dashed, draw=danger!70] (toxi) -- (mysql);
  \draw[arrow, dashed, draw=danger!70] (toxi) -- (redis);
  \draw[arrow, dashed, draw=danger!70] (toxi) -- (kafka);
\end{tikzpicture}
\caption{High-level three-layer architecture. Dashed red arrows represent
  Toxiproxy fault injection paths.}
\label{fig:hla}
\end{figure}

\subsection{Container Map}

The system runs as \textbf{10 Docker containers} on a custom bridge
network named \texttt{stock-network} (Table~\ref{tab:containers}).

\begin{table}[H]
\centering
\caption{Docker Compose service inventory}
\label{tab:containers}
\begin{tabularx}{\textwidth}{@{}lllX@{}}
\toprule
\textbf{Container}       & \textbf{Image}              & \textbf{Port(s)}  & \textbf{Role} \\
\midrule
\texttt{kafka-kraft}     & apache/kafka:4.1.1          & 9092              & Event streaming (KRaft, no ZooKeeper) \\
\texttt{mysql-db}        & mysql:8.0                   & 3306              & Source-of-truth + Outbox table \\
\texttt{redis-master}    & redis:7-alpine              & 6379              & Cache + Sensor data store (R/W) \\
\texttt{redis-replica}   & redis:7-alpine              & 6380              & Read-only standby \\
\texttt{toxiproxy}       & shopify/toxiproxy:latest    & 8474, 3307, 26379, 29093 & Network partition injection \\
\texttt{iot-simulator}   & custom Python               & —                 & Warehouse temperature generator \\
\texttt{stock-app}       & custom Spring Boot          & 8080              & Business logic + REST API + SSE \\
\texttt{prometheus}      & prom/prometheus:latest      & 9090              & Metrics scraping (15 s interval) \\
\texttt{grafana}         & grafana/grafana:latest      & 3001              & Dashboards (Stock + IoT) \\
\texttt{stock-frontend}  & custom Nginx + React        & 3000              & SPA + reverse proxy \\
\bottomrule
\end{tabularx}
\end{table}

%% ============================================================
\section{Component Architecture}
\label{sec:components}
%% ============================================================

\subsection{Backend: Spring Boot Application}

The backend is built on \tech{Spring Boot 3.x} (Java 17, Gradle 8.5)
and is structured into four horizontal layers: \emph{Controller},
\emph{Service}, \emph{Repository}, and \emph{Configuration}
(Figure~\ref{fig:backend}).

\begin{figure}[H]
\centering
\begin{tikzpicture}[node distance=0.4cm and 0.3cm, scale=0.82, every node/.style={scale=0.82}]

  %% Controller row
  \node[svcbox=secondary, minimum width=2.6cm] (sc)  at (-3.9, 0) {StockCtrl};
  \node[svcbox=secondary, minimum width=2.6cm] (sec) at (-1.1, 0) {SensorCtrl\\(SSE)};
  \node[svcbox=secondary, minimum width=2.6cm] (lc)  at (1.7, 0)  {LoadTest\\Ctrl};
  \node[svcbox=secondary, minimum width=2.6cm] (rl)  at (4.5, 0)  {RateLimit\\Interceptor};
  \node[layerlabel] at (-6.5,0) {Controllers};

  %% Service row
  \node[svcbox=primary, minimum width=2.6cm] (ss)  at (-3.9,-1.8) {StockService};
  \node[svcbox=primary, minimum width=2.6cm] (cs)  at (-1.1,-1.8) {CacheService\\(Lettuce)};
  \node[svcbox=primary, minimum width=2.6cm] (kps) at (1.7,-1.8)  {KafkaProducer\\Service};
  \node[svcbox=primary, minimum width=2.6cm] (wml) at (4.5,-1.8)  {Warehouse\\MetricsListener};
  \node[layerlabel] at (-6.5,-1.8) {Services};

  %% Outbox & Rate
  \node[svcbox=accent, minimum width=2.6cm] (obs) at (-1.1,-3.4) {Outbox\\Scheduler};
  \node[svcbox=accent, minimum width=2.6cm] (rls) at (1.7,-3.4)  {RateLimit\\Service};
  \node[layerlabel] at (-6.5,-3.4) {Schedulers};

  %% Infrastructure row
  \node[svcbox=success, minimum width=2cm] (mdb) at (-3.9,-5.0) {MySQL};
  \node[svcbox=warning, minimum width=2cm] (rdb) at (-1.3,-5.0) {Redis};
  \node[svcbox=danger,  minimum width=2cm] (kbr) at (1.1,-5.0)  {Kafka};
  \node[svcbox=neutral, minimum width=2cm] (prd) at (3.5,-5.0)  {Prometheus};
  \node[layerlabel] at (-6.5,-5.0) {Infrastructure};

  %% Arrows — controllers → services
  \draw[arrow] (sc)  -- (ss);
  \draw[arrow] (sec) -- (cs);
  \draw[arrow] (lc)  -- (kps);

  %% services → infra
  \draw[arrow] (ss)  -- (mdb);
  \draw[arrow] (cs)  -- (rdb);
  \draw[arrow] (kps) -- (kbr);
  \draw[arrow] (obs) -- (kbr);
  \draw[arrow] (wml) -- (rdb);

  %% Outbox
  \draw[arrow] (ss) -- (obs);
\end{tikzpicture}
\caption{Spring Boot application internal component diagram.}
\label{fig:backend}
\end{figure}

\subsubsection{REST API Endpoints}

\begin{lstlisting}[style=javaStyle, caption={Controller endpoint overview}]
// StockController
GET    /api/stocks          -> list all stocks
GET    /api/stocks/{id}     -> get by id
POST   /api/stocks          -> create stock
PUT    /api/stocks/{id}     -> update stock
DELETE /api/stocks/{id}     -> delete stock

// SensorController (SSE + REST)
GET    /api/sensors/stream/{deviceId}       -> SSE real-time stream
GET    /api/sensors/current/{deviceId}      -> latest Redis reading
GET    /api/sensors/alerts/{deviceId}       -> alert list
GET    /api/sensors/alerts/{deviceId}/clear -> clear alerts
GET    /api/sensors/health/{deviceId}       -> sensor health

// LoadTestController
POST   /api/load-test/kafka?count=10000     -> throughput benchmark
POST   /api/load-test/kafka/batch           -> batch benchmark
\end{lstlisting}

\subsubsection{Key Dependencies}

\begin{table}[H]
\centering
\caption{Spring Boot key dependencies}
\label{tab:deps}
\begin{tabularx}{\textwidth}{@{}llX@{}}
\toprule
\textbf{Artifact}                         & \textbf{Purpose}              & \textbf{Notes} \\
\midrule
spring-boot-starter-web                   & REST API                      & Embedded Tomcat \\
spring-boot-starter-data-jpa              & ORM + Transactions            & HikariCP (10 conn.) \\
spring-kafka                              & Event streaming               & Producer + Consumer \\
spring-boot-starter-data-redis            & Cache layer                   & Lettuce async client \\
spring-boot-starter-actuator              & Health + metrics              & /actuator/prometheus \\
micrometer-registry-prometheus            & Metrics export                & 15 s scrape interval \\
bucket4j-core:8.10.1                      & Rate limiting                 & Token-bucket 50 req/s \\
mysql-connector-j                         & JDBC driver                   & MySQL 8 \\
\bottomrule
\end{tabularx}
\end{table}

\subsection{Data Layer}

\subsubsection{MySQL — Source of Truth}

MySQL 8.0 serves as the \emph{source of truth} with ACID-compliant
transactions.
Two schemas are maintained: the primary \texttt{stocks} table and an
\texttt{outbox\_events} table that implements the Transactional Outbox
Pattern (Section~\ref{sec:outbox}).

\begin{lstlisting}[style=sqlStyle, caption={Database schema}]
CREATE TABLE stocks (
    id           BIGINT AUTO_INCREMENT PRIMARY KEY,
    product_name VARCHAR(255) UNIQUE NOT NULL,
    quantity     INT NOT NULL,
    price        DECIMAL(19,2) NOT NULL,
    created_at   TIMESTAMP DEFAULT CURRENT_TIMESTAMP,
    updated_at   TIMESTAMP DEFAULT CURRENT_TIMESTAMP
                   ON UPDATE CURRENT_TIMESTAMP
);

CREATE TABLE outbox_events (
    id          BIGINT AUTO_INCREMENT PRIMARY KEY,
    event_id    VARCHAR(36) UNIQUE NOT NULL,
    event_type  VARCHAR(20)  NOT NULL,           -- CREATED|UPDATED|DELETED
    stock_id    BIGINT NOT NULL,
    payload     TEXT NOT NULL,                   -- JSON
    status      ENUM('PENDING','SENT','FAILED')
                  NOT NULL DEFAULT 'PENDING',
    retry_count INT NOT NULL DEFAULT 0,
    last_error  VARCHAR(500),
    created_at  TIMESTAMP NOT NULL,
    processed_at TIMESTAMP
);
\end{lstlisting}

\subsubsection{Redis — Cache Layer}

Redis 7 runs in a Master/Replica topology.
The master handles all reads and writes; the replica provides a
read-only standby.
Three caching strategies are employed depending on the operational
context (Figure~\ref{fig:cache}).

\begin{figure}[H]
\centering
\begin{tikzpicture}[node distance=1.0cm, scale=0.84, every node/.style={scale=0.84}]

  \node[svcbox=primary, minimum width=3cm]  (svc)  at (0,0)   {StockService};
  \node[svcbox=accent,  minimum width=2.4cm](redis) at (4,0)   {Redis Cache};
  \node[svcbox=success, minimum width=2.4cm](mysql) at (4,-2)  {MySQL DB};

  %% Write-Through
  \draw[arrow, draw=success!80] (svc.east) to[bend left=20]
    node[above, font=\tiny\color{success!80}]{Write-Through (Create/Update)} (redis.west);

  %% Read hit
  \draw[arrow, draw=accent!80] (redis.west) to[bend left=20]
    node[below, font=\tiny\color{accent!80}]{Cache HIT → return} (svc.east);

  %% Read miss chain
  \draw[arrow] (svc) -- node[left,font=\tiny]{MISS} (mysql |- svc) -- (mysql);
  \draw[arrow, draw=neutral!60] (mysql) to[bend right=20]
    node[right,font=\tiny]{populate} (redis);

  %% Fallback
  \node[font=\tiny\itshape, text=danger, align=center] at (2,-3.0)
    {Cache unavailable → direct DB fallback (no cache write)};
  \draw[-Stealth, dashed, draw=danger!70] (svc) -- (2,-2.8);

\end{tikzpicture}
\caption{Redis caching strategies: Write-Through on creation, Read-Through on
  cache hit, Cache-Aside fallback during partition.}
\label{fig:cache}
\end{figure}

Cache TTL is set to \textbf{300 seconds}. The Lettuce client is configured
with a 5-second connection timeout to trigger graceful degradation
before overwhelming downstream services.

\subsubsection{Kafka — Event Bus (KRaft Mode)}

Kafka 4.1.1 operates in \emph{KRaft mode} — no ZooKeeper dependency.
Two topics carry the system's event traffic (Table~\ref{tab:topics}).

\begin{table}[H]
\centering
\caption{Kafka topics and configurations}
\label{tab:topics}
\begin{tabularx}{\textwidth}{@{}llXl@{}}
\toprule
\textbf{Topic}         & \textbf{Partitions} & \textbf{Producer} & \textbf{Consumer} \\
\midrule
\texttt{stock-events}      & 3 & Outbox Scheduler (Java, sync) & — (external) \\
\texttt{warehouse-metrics} & 3 & IoT Simulator (Python, JSON)  & WarehouseMetricsListener \\
\bottomrule
\end{tabularx}
\end{table}

Producer configuration enforces maximum durability:
\texttt{acks=all}, \texttt{enable.idempotence=true},
\texttt{retries=3}, \texttt{linger.ms=10}.

\subsection{Frontend: React + Nginx}

The frontend is a \tech{React 18} SPA bundled with \tech{Vite 6}.
It is served by an \tech{Nginx} reverse proxy that:
\begin{enumerate}[noitemsep]
  \item Forwards \texttt{/api/} to the Spring Boot backend (\texttt{app:8080}),
        eliminating CORS issues entirely in Docker.
  \item Disables proxy buffering for \textbf{Server-Sent Events (SSE)}
        to deliver real-time sensor data to the browser.
  \item Sets \texttt{proxy\_read\_timeout 300s} to keep SSE connections alive.
\end{enumerate}

The UI provides three views:
\begin{itemize}[noitemsep]
  \item \textbf{Stock Management} — CRUD operations with real-time list refresh.
  \item \textbf{IoT Sensor Dashboard} — Live SSE temperature stream, alert
        management, and an HTML5 Canvas time-series chart (last 60 readings).
  \item \textbf{Metrics Dashboard} — Embedded Grafana iframes with a
        toggle between stock and IoT sensor dashboards.
\end{itemize}

\subsection{Fault Injection: Toxiproxy}

\tech{Toxiproxy} acts as a transparent TCP proxy in front of MySQL,
Redis, and Kafka (Figure~\ref{fig:toxi}).
Three toxic types are used to simulate distinct network failure modes:

\begin{table}[H]
\centering
\caption{Toxiproxy toxic types used in experiments}
\label{tab:toxics}
\begin{tabularx}{\textwidth}{@{}llX@{}}
\toprule
\textbf{Toxic}     & \textbf{Effect} & \textbf{Experiment Purpose} \\
\midrule
\texttt{timeout}   & Drop all packets (timeout=0) & Complete component partition \\
\texttt{latency}   & Add round-trip delay + jitter & Throughput degradation study \\
\texttt{bandwidth} & Throttle to 1 KB/s & Slow-link simulation \\
\bottomrule
\end{tabularx}
\end{table}

\begin{figure}[H]
\centering
\begin{tikzpicture}[node distance=0.5cm and 0.3cm, scale=0.84, every node/.style={scale=0.84}]

  \node[svcbox=primary, minimum width=3cm] (app) at (0,0) {Spring Boot};

  \node[svcbox=danger, minimum width=2cm] (tm) at (-3,-2)
    {Toxiproxy\\:3307};
  \node[svcbox=danger, minimum width=2cm] (tr) at (0,-2)
    {Toxiproxy\\:26379};
  \node[svcbox=danger, minimum width=2cm] (tk) at (3,-2)
    {Toxiproxy\\:29093};

  \node[svcbox=success, minimum width=2cm] (mysql) at (-3,-4) {MySQL};
  \node[svcbox=accent,  minimum width=2cm] (redis) at (0,-4)  {Redis};
  \node[svcbox=warning, minimum width=2cm] (kafka) at (3,-4)  {Kafka};

  \draw[arrow] (app) -- (tm);
  \draw[arrow] (app) -- (tr);
  \draw[arrow] (app) -- (tk);
  \draw[arrow] (tm)  -- node[right,font=\tiny]{JDBC}    (mysql);
  \draw[arrow] (tr)  -- node[right,font=\tiny]{Lettuce} (redis);
  \draw[arrow] (tk)  -- node[right,font=\tiny]{Kafka}   (kafka);

  \node[font=\tiny, text=danger, align=center] at (0,-5)
    {Toxics injected via REST API: POST /proxies/\{name\}/toxics};
\end{tikzpicture}
\caption{Toxiproxy topology. All infrastructure connections from the
  application pass through a proxy for controllable fault injection.}
\label{fig:toxi}
\end{figure}

%% ============================================================
\section{Data Flow Patterns}
\label{sec:dataflow}
%% ============================================================

\subsection{Stock Creation — Happy Path}

Figure~\ref{fig:create} traces the sequence of operations when a client
creates a new stock item.

\begin{figure}[H]
\centering
\begin{tikzpicture}[node distance=0.5cm, scale=0.84, every node/.style={scale=0.84}]
  \foreach \x/\lbl in {0/Client, 2.4/StockCtrl, 4.8/StockService,
                        7.2/MySQL, 9.6/Redis, 12/Kafka}{
    \node[svcbox=neutral, minimum width=2cm, font=\small\bfseries]
      (\lbl) at (\x,0) {\lbl};
    \draw[dashed, draw=neutral!40] (\x,-0.4) -- (\x,-5.6);
  }

  %% Messages
  \draw[arrow] (0,-0.9)  -- node[above,font=\tiny]{POST /api/stocks} (2.4,-0.9);
  \draw[arrow] (2.4,-1.4)-- node[above,font=\tiny]{createStock()}    (4.8,-1.4);
  \draw[arrow] (4.8,-2.0)-- node[above,font=\tiny]{save(stock) TX}   (7.2,-2.0);
  \draw[arrow] (7.2,-2.5)-- node[above,font=\tiny\color{success}]{commit ✓} (4.8,-2.5);
  \draw[arrow] (4.8,-3.1)-- node[above,font=\tiny]{SET stock:id TTL300}(9.6,-3.1);
  \draw[arrow] (9.6,-3.6)-- node[above,font=\tiny\color{success}]{OK ✓} (4.8,-3.6);
  \draw[arrow] (4.8,-4.2)-- node[above,font=\tiny]{sendEvent() async} (12,-4.2);
  \draw[arrow] (4.8,-5.0)-- node[above,font=\tiny\color{success}]{201 Created ✓}(0,-5.0);
\end{tikzpicture}
\caption{Stock creation sequence: synchronous DB write, synchronous cache
  population, and asynchronous Kafka event dispatch.}
\label{fig:create}
\end{figure}

\subsection{Stock Retrieval — Cache Hit vs. Cache Miss}

On a \textbf{cache hit}, the response latency is under 10 ms because
the data is returned directly from Redis without touching MySQL.
On a \textbf{cache miss}, the service falls back to MySQL (~50 ms),
then populates the cache for future requests (write-through).

\subsection{Partition Handling — Redis Down}

When Toxiproxy drops all Redis packets:

\begin{enumerate}[noitemsep]
  \item CacheService attempts \texttt{GET stock:\{id\}}.
  \item Lettuce connection times out after 5000 ms.
  \item A \texttt{RedisConnectionException} is caught silently.
  \item The \texttt{redis.reconnect.total} counter is incremented.
  \item MySQL is queried directly — response is returned to client.
  \item \emph{No cache write} occurs during the partition to prevent
        stale-entry pollution.
\end{enumerate}

Result: \textbf{full functional availability} at the cost of higher
latency (~5050 ms per request).

%% ============================================================
\section{CAP Theorem Analysis}
\label{sec:cap}
%% ============================================================

\subsection{Component-Level CAP Positions}

Figure~\ref{fig:cap} illustrates where each major component sits in the
CAP triangle.

\begin{figure}[H]
\centering
\begin{tikzpicture}[scale=1.0]
  %% Triangle
  \coordinate (C) at (0, 3.5);
  \coordinate (A) at (-3, 0);
  \coordinate (P) at (3, 0);
  \draw[thick, draw=neutral!60] (C) -- (A) -- (P) -- cycle;

  %% Vertex labels
  \node[above, font=\bfseries\large, text=primary] at (C) {C};
  \node[below left,  font=\bfseries\large, text=success] at (A) {A};
  \node[below right, font=\bfseries\large, text=warning] at (P) {P};
  \node[above, font=\small, text=primary]   at ($(C)+(0,0.2)$)  {Consistency};
  \node[below, font=\small, text=success]   at ($(A)+(0,-0.2)$) {Availability};
  \node[below, font=\small, text=warning]   at ($(P)+(0,-0.2)$) {Part. Tolerance};

  %% CP midpoint → MySQL
  \coordinate (CP) at ($(C)!0.45!(P)$);
  \fill[success!80] (CP) circle (5pt);
  \node[right, font=\small\bfseries, text=success!80] at ($(CP)+(0.1,0)$) {MySQL (CP)};

  %% AP midpoint → Redis
  \coordinate (AP) at ($(A)!0.45!(P)$);
  \fill[accent!80] (AP) circle (5pt);
  \node[below, font=\small\bfseries, text=accent!80] at ($(AP)+(0,-0.1)$) {Redis (AP)};

  %% System overall → slightly toward AP
  \coordinate (SYS) at (-1.0, 0.9);
  \fill[primary] (SYS) circle (6pt);
  \node[left, font=\small\bfseries, text=primary] at ($(SYS)+(-0.1,0)$) {System};

  %% Kafka note
  \node[font=\tiny\itshape, text=warning, align=center] at (0,-0.7)
    {Kafka: CP (native) → AP via graceful degradation};
\end{tikzpicture}
\caption{CAP triangle showing positions of individual components and the
  overall system (AP-prioritised).}
\label{fig:cap}
\end{figure}

\begin{table}[H]
\centering
\caption{Component CAP characteristics summary}
\label{tab:cap}
\begin{tabularx}{\textwidth}{@{}llXl@{}}
\toprule
\textbf{Component} & \textbf{CAP} & \textbf{Partition Behaviour} & \textbf{Strategy} \\
\midrule
MySQL  & CP & Fail → HTTP 500, no data access & Fail-fast \\
Redis  & AP & Fail → DB fallback, degraded perf. & Cache-Aside fallback \\
Kafka  & CP$\to$AP & Fail → event loss, stock op. succeeds & Fire-and-forget + log \\
\bottomrule
\end{tabularx}
\end{table}

\subsection{System-Level Decision: AP with CP Fallback}

The system is designed to maximise availability:
\begin{itemize}[noitemsep]
  \item Stock CRUD operations continue even if Redis \emph{and} Kafka are down.
  \item Only a MySQL partition causes a hard failure (HTTP 500), because
        the database is the sole source of truth and there is no safe
        fallback for persistence.
  \item Eventual consistency is preferred over blocking writes.
\end{itemize}

%% ============================================================
\section{Observability Stack}
\label{sec:observability}
%% ============================================================

\subsection{Prometheus Metrics}

Spring Boot Actuator exposes a \texttt{/actuator/prometheus} endpoint.
Prometheus scrapes it every 15 seconds.
Beyond JVM and HTTP auto-instrumentation, the system defines 8 custom
Micrometer metrics (Table~\ref{tab:metrics}).

\begin{table}[H]
\centering
\caption{Custom Micrometer metrics}
\label{tab:metrics}
\begin{tabularx}{\textwidth}{@{}llX@{}}
\toprule
\textbf{Metric Name}              & \textbf{Type}   & \textbf{Description} \\
\midrule
\texttt{cache.hit.miss.total}     & Counter  & Cache hits/misses tagged by result \\
\texttt{cache.operation.duration} & Timer    & Redis op latency (get/set/delete) \\
\texttt{redis.reconnect.total}    & Counter  & Lettuce reconnection attempts \\
\texttt{kafka.producer.send.total}& Counter  & Events sent, tagged success/failure \\
\texttt{stock.create.duration}    & Timer    & End-to-end stock creation latency \\
\texttt{stock.update.duration}    & Timer    & End-to-end stock update latency \\
\texttt{outbox.events.total}      & Counter  & Outbox events tagged sent/failed/retried \\
\texttt{outbox.pending.gauge}     & Gauge    & Current pending Outbox event count \\
\bottomrule
\end{tabularx}
\end{table}

\subsection{Grafana Dashboards}

Two dashboards are provisioned automatically via Grafana's provisioning
API:

\begin{enumerate}
  \item \textbf{Stock Metrics Dashboard} — Six stat panels (App health,
        cache hit rate, Kafka success rate, throughput, Redis reconnects,
        Kafka failures) plus time-series graphs for HTTP rate, cache
        latency, cache hit/miss stacked area, Kafka throughput, and
        JVM heap.

  \item \textbf{IoT Sensor Metrics Dashboard} — Temperature time series,
        alert rate, partition health, and consumer lag gauges.
\end{enumerate}

Alert thresholds:
\begin{itemize}[noitemsep]
  \item Cache hit rate $< 50\%$ → \textcolor{warning}{\textbf{Warning}}
  \item Kafka success rate $< 80\%$ → \textcolor{danger}{\textbf{Critical}}
  \item Redis reconnects $> 0$ → \textcolor{danger}{\textbf{Critical}}
\end{itemize}

%% ============================================================
\section{IoT Sensor Pipeline}
\label{sec:iot}
%% ============================================================

Figure~\ref{fig:iot} illustrates the end-to-end path of IoT sensor data.

\begin{figure}[H]
\centering
\begin{tikzpicture}[node distance=0.3cm and 0.2cm, scale=0.84, every node/.style={scale=0.84}]

  \node[svcbox=warning, minimum width=3cm] (sim)
    at (0,0) {Python IoT\\Simulator};
  \node[svcbox=warning, minimum width=3cm] (kf)
    at (4,0) {Kafka\\warehouse-metrics};
  \node[svcbox=primary, minimum width=3cm] (wml)
    at (8,0)  {WarehouseMetrics\\Listener};
  \node[svcbox=accent,  minimum width=3cm] (rd)
    at (8,-2) {Redis\\sensor:\{id\}\\alert:\{id\}};
  \node[svcbox=secondary,minimum width=3cm] (ctl)
    at (4,-2) {SensorController\\SSE stream};
  \node[svcbox=neutral, minimum width=3cm] (fe)
    at (0,-2) {React\\Dashboard};

  \draw[arrow] (sim) -- node[above,font=\tiny]{JSON 1 s} (kf);
  \draw[arrow] (kf)  -- node[above,font=\tiny]{consume} (wml);
  \draw[arrow] (wml) -- node[right,font=\tiny]{SET / LPUSH} (rd);
  \draw[arrow] (rd)  -- node[above,font=\tiny]{LRANGE alerts} (ctl);
  \draw[arrow] (ctl) -- node[above,font=\tiny]{SSE events} (fe);

  \node[msgbox, below=0.5cm of sim, align=left] {
    \{"cihazId":"depo-sensor-1",\\
    \,"sicaklik":24.5,\\
    \,"zaman":"2026-04-09T12:03:10Z"\}
  };
\end{tikzpicture}
\caption{IoT sensor data pipeline from Python simulator to React dashboard
  via Kafka, Redis, and Server-Sent Events.}
\label{fig:iot}
\end{figure}

\subsubsection{Python Sensor Simulator}

The simulator (\texttt{sensor\_simulator.py}) runs in its own Docker
container and publishes a temperature reading every \textbf{1 second}
to the \texttt{warehouse-metrics} Kafka topic.

\begin{lstlisting}[style=javaStyle, language=Python,
  caption={Core simulation loop (Python)}]
def generate_sensor_data():
    temperature = round(random.uniform(15.0, 30.0), 1)
    return {
        "cihazId":   DEVICE_ID,          # "depo-sensor-1"
        "sicaklik":  temperature,        # Celsius
        "zaman":     datetime.utcnow()   # ISO-8601 + Z
    }

while True:
    send_to_kafka(producer, generate_sensor_data())
    time.sleep(1)
\end{lstlisting}

\subsubsection{Alert Generation}

The \texttt{WarehouseMetricsListener} Kafka consumer triggers an alert
whenever the temperature exceeds \textbf{30 °C}:

\begin{lstlisting}[style=javaStyle,
  caption={Temperature alert logic in Java}]
if (dto.getSicaklik() > 30.0) {
    redisTemplate.opsForList()
        .leftPush("alert:" + dto.getCihazId(),
                  "HIGH TEMP: " + dto.getSicaklik() + " C at " + dto.getZaman());
}
\end{lstlisting}

Alerts are stored as a Redis List and exposed via the SensorController
REST endpoint, which the React dashboard polls.

%% ============================================================
\section{Resilience Patterns}
\label{sec:resilience}
%% ============================================================

\subsection{Transactional Outbox Pattern}
\label{sec:outbox}

The Outbox Pattern (Figure~\ref{fig:outbox}) guarantees \textbf{at-least-once}
delivery of stock domain events to Kafka, even during broker outages.

\begin{figure}[H]
\centering
\begin{tikzpicture}[node distance=0.4cm, scale=0.84, every node/.style={scale=0.84}]

  \node[svcbox=primary, minimum width=3cm]  (ss) at (0,0)    {StockService};
  \node[svcbox=success, minimum width=2.4cm](db) at (3.5,0)  {MySQL TX};
  \node[svcbox=accent,  minimum width=2.8cm](ob) at (3.5,-1.8){outbox\_events\\PENDING};
  \node[svcbox=primary, minimum width=3cm]  (sc) at (0,-3.6) {OutboxScheduler\\(every 5 s)};
  \node[svcbox=warning, minimum width=2.4cm](kf) at (3.5,-3.6){Kafka\\stock-events};
  \node[svcbox=neutral, minimum width=2.8cm](st) at (3.5,-5.2){status = SENT\\or FAILED (>5)};

  \draw[arrow] (ss) -- node[above,font=\tiny]{save(stock)} (db);
  \draw[arrow] (db) -- node[right,font=\tiny]{\textregistered same TX} (ob);
  \draw[arrow] (sc) -- node[above,font=\tiny]{SELECT PENDING LIMIT 50} (ob |- sc) -- (ob);
  \draw[arrow] (sc) -- node[above,font=\tiny]{sendSync()} (kf);
  \draw[arrow] (kf) -- node[right,font=\tiny]{ACK} (st);
  \draw[arrow, dashed, draw=danger!70] (kf.south) -- node[right,font=\tiny\color{danger}]{retryCount++} (st.north);
\end{tikzpicture}
\caption{Transactional Outbox Pattern: stock and event writes share one
  database transaction; a background scheduler retries delivery up to
  five times before marking an event as FAILED.}
\label{fig:outbox}
\end{figure}

\subsection{Rate Limiting — Token Bucket}

A Bucket4j-based token-bucket rate limiter enforces a global limit of
\textbf{50 requests/second}, configured via application properties:

\begin{lstlisting}[style=yamlStyle, caption={Rate limiter configuration}]
rate-limit.global:
  capacity:            50   # Max bucket size
  refill-tokens:       50   # Tokens restored per interval
  refill-duration-ms: 1000  # Refill interval (1 s)
\end{lstlisting}

Excess requests receive HTTP \texttt{429 Too Many Requests}.

%% ============================================================
\section{Testing \& Benchmarking}
\label{sec:testing}
%% ============================================================

\subsection{Fault-Injection Experiments}

Three partition scenarios were systematically tested:

\paragraph{Redis Partition.}
Injecting a \texttt{timeout} toxic with \texttt{timeout=0} on the
Redis proxy causes all cache operations to time out after 5 000 ms.
The system falls back to MySQL transparently. Cache-miss metrics spike
in Grafana; average GET latency increases from 10 ms to~5 050 ms.
Upon removing the toxic, the cache warms organically as new reads
populate entries, and latency normalises within seconds.

\paragraph{Kafka Partition.}
With a Kafka timeout toxic active, the Outbox Scheduler accumulates
\texttt{PENDING} events in MySQL while retrying every 5 seconds.
\texttt{stock.create} and \texttt{stock.update} operations continue to
return HTTP 201/200 — clients are unaffected.
Once the toxic is removed, the Scheduler delivers all backlogged events;
zero event loss is observed.

\paragraph{MySQL Partition.}
A MySQL timeout toxic causes all JDBC calls to throw
\texttt{SQLException}.
The \texttt{@Transactional} annotation triggers a rollback; clients
receive HTTP 500.
This confirms the CP behaviour of MySQL: \emph{no availability, perfect
consistency}. Recovery is instantaneous upon toxic removal.

\subsection{Throughput Benchmarks}

Load tests were performed by sending 10 000 Kafka messages through the
load-test endpoint (Table~\ref{tab:bench}).

\begin{table}[H]
\centering
\caption{Kafka throughput under varying network conditions}
\label{tab:bench}
\begin{tabularx}{\textwidth}{@{}lrrrl@{}}
\toprule
\textbf{Scenario}          & \textbf{Throughput} & \textbf{Duration} & \textbf{Success Rate} & \textbf{Notes} \\
\midrule
Normal (no toxic)          & 200 ops/s & 2–3 s   & 100\% & Baseline \\
Kafka 500 ms latency       & 125 ops/s & ~8 s    & 100\% & 62.5\% of baseline \\
Kafka 2000 ms latency      &  25 ops/s & ~40 s   & 100\% & 12.5\% of baseline \\
Kafka complete partition    &   0 ops/s & N/A     &   0\% & Events queued in Outbox \\
\bottomrule
\end{tabularx}
\end{table}

The linear relationship $\text{throughput} \propto 1/\text{latency}$
holds under latency injection, confirming the synchronous send semantics
of the Outbox Scheduler.

\subsection{Latency Measurements}

\begin{table}[H]
\centering
\caption{HTTP endpoint response times under normal and partition conditions}
\label{tab:latency}
\begin{tabular}{@{}lrrl@{}}
\toprule
\textbf{Endpoint}            & \textbf{Cache HIT} & \textbf{Cache MISS} & \textbf{Redis Down} \\
\midrule
GET /api/stocks/\{id\}       & $<$10 ms  & $\sim$50 ms  & $\sim$5050 ms \\
POST /api/stocks             & $\sim$80 ms & $\sim$80 ms & $\sim$85 ms \\
GET /api/sensors/stream/...  & SSE (persistent) & — & — \\
\bottomrule
\end{tabular}
\end{table}

\subsection{Resource Utilisation}

JVM heap usage stabilises at approximately \textbf{40 MB} (2\% of the
2 048 MB maximum) at 200 ops/s, demonstrating the efficiency of the
Lettuce async client and non-blocking Kafka producer.
Connection pools are sized conservatively: 10 MySQL connections
(HikariCP), 8 Redis connections (Lettuce), 5 in-flight Kafka producer
requests.

%% ============================================================
\section{Key Design Decisions \& Trade-offs}
\label{sec:design}
%% ============================================================

\subsection{Kafka \texttt{acks=all} with Graceful Degradation}

Using \texttt{acks=all} maximises durability when Kafka is healthy —
the producer waits for the leader and all in-sync replicas to
acknowledge.
When Kafka becomes unavailable, the exception handler increments the
failure counter and logs the error \emph{without rethrowing}, preserving
stock-operation availability.
This effectively uses a CP technology in an AP manner.

\subsection{Hybrid Redis Strategy: Write-Through + Cache-Aside}

On creation and update, the cache is synchronously populated
immediately after the database write (\emph{write-through}), ensuring
subsequent reads are cache-hits.
During a Redis partition, the \emph{cache-aside} branch skips the
cache-write to avoid storing potentially stale entries whose TTL would
not begin expiring until recovered.

\subsection{KRaft Mode — No ZooKeeper}

Kafka 4.1.1 uses the built-in Raft consensus mechanism (\emph{KRaft})
instead of an external ZooKeeper ensemble.
This simplifies the compose file by eliminating a separate ZooKeeper
container, reduces startup time, and reflects the production trajectory
of the Kafka project.

\subsection{Single Kafka Broker}

A single broker is used for development and experiment simplicity.
Toxiproxy compensates for the absence of multi-broker fault tolerance
by providing controllable network partitions.
A production deployment would use a 3-broker cluster with replication
factor~3.

\subsection{React + Native Canvas Charting}

Third-party charting libraries (Chart.js, Recharts) were deliberately
avoided.
The temperature time-series chart is rendered on an HTML5 Canvas,
eliminating dependencies, reducing bundle size, and maintaining full
control over the rendering pipeline required for 1 Hz SSE update rates.

%% ============================================================
\section{CI/CD Pipeline}
\label{sec:cicd}
%% ============================================================

A GitHub Actions workflow (\texttt{.github/workflows/ci.yml}) runs on
every push and pull request to \texttt{master}:

\begin{enumerate}[noitemsep]
  \item \textbf{Gradle Build} — compiles Java sources and runs JUnit tests
        with Spring Kafka Test embedded broker.
  \item \textbf{Docker Build} — builds the Spring Boot image to validate
        the \texttt{Dockerfile}.
  \item \textbf{Integration Smoke Test} — starts the full Docker Compose
        stack and sends a health-check HTTP request to verify the application
        starts correctly.
\end{enumerate}

%% ============================================================
\section{Future Work}
\label{sec:future}
%% ============================================================

\begin{enumerate}[leftmargin=1.8cm, label=\textbf{F\arabic*.}]
  \item \textbf{Redis Sentinel} — automatic failover from master to
        standby; eliminate the manual promotion step.
  \item \textbf{Multi-Broker Kafka} — 3-broker cluster with replication
        factor 3 for genuine Kafka-level fault tolerance.
  \item \textbf{Circuit Breakers (Resilience4j)} — prevent cascading
        failures when Redis or MySQL degrade slowly rather than failing
        instantly.
  \item \textbf{Dead-Letter Queue} — persist permanently-failed Outbox
        events for manual replay and auditing.
  \item \textbf{Distributed Tracing (Zipkin/Jaeger)} — correlate spans
        across HTTP, Kafka, and Redis to trace latency sources end-to-end.
  \item \textbf{Reduced Redis Timeout} — lower the Lettuce connection
        timeout from 5 s to 2 s to shorten the degraded-mode response
        window.
  \item \textbf{Multi-Tenant IoT Sensors} — support multiple \texttt{deviceId}
        streams simultaneously; partition \texttt{warehouse-metrics} by
        device.
\end{enumerate}

%% ============================================================
\section{Conclusion}
\label{sec:conclusion}
%% ============================================================

This paper has documented a pragmatic, \textbf{AP-focused} distributed
system that demonstrates CAP Theorem trade-offs through live fault
injection.

The principal contributions are:

\begin{itemize}
  \item A ten-container Docker Compose architecture covering the full
        technology stack of a modern microservice system.
  \item Systematic Toxiproxy experiments that make CAP trade-offs
        \emph{measurable} rather than theoretical.
  \item A \textbf{Transactional Outbox Pattern} that upgrades stock-event
        delivery from \emph{at-most-once} to \emph{at-least-once} without
        sacrificing availability.
  \item An end-to-end IoT pipeline (Python → Kafka → Redis → SSE → React)
        that demonstrates heterogeneous producer/consumer integration.
  \item 30+ Prometheus metrics and two Grafana dashboards that surface
        partition effects in sub-15-second resolution.
\end{itemize}

The system remains \textbf{fully functional} under Redis or Kafka
partitions — the two failure modes most relevant to an AP-prioritised
architecture — while correctly failing fast under MySQL partitions where
no safe consistency-preserving fallback exists.

%% ============================================================
%%  References
%% ============================================================
\newpage
\begin{thebibliography}{9}

\bibitem{brewer2000towards}
E.~A. Brewer,
\newblock ``Towards robust distributed systems,''
\newblock in \emph{Proc. 19th Annual ACM Symposium on Principles of Distributed
  Computing (PODC)}, Portland, OR, 2000.

\bibitem{gilbert2002brewer}
S.~Gilbert and N.~Lynch,
\newblock ``Brewer's conjecture and the feasibility of consistent, available,
  partition-tolerant web services,''
\newblock \emph{ACM SIGACT News}, vol.~33, no.~2, pp.~51--59, 2002.

\bibitem{kleppmann2017ddia}
M.~Kleppmann,
\newblock \emph{Designing Data-Intensive Applications}.
\newblock O'Reilly Media, 2017.

\bibitem{springsource}
VMware,
\newblock \emph{Spring Boot Reference Documentation 3.x},
\newblock \url{https://docs.spring.io/spring-boot/}, 2024.

\bibitem{kafka}
Apache Software Foundation,
\newblock \emph{Apache Kafka Documentation 4.1},
\newblock \url{https://kafka.apache.org/documentation/}, 2024.

\bibitem{toxiproxy}
Shopify Inc.,
\newblock \emph{Toxiproxy — A TCP proxy to simulate network and system
  conditions},
\newblock \url{https://github.com/Shopify/toxiproxy}, 2023.

\bibitem{richardson2018microservices}
C.~Richardson,
\newblock \emph{Microservices Patterns}.
\newblock Manning Publications, 2018.

\end{thebibliography}

%% ============================================================
\appendix
\section{Technology Version Matrix}
\label{app:versions}
%% ============================================================

\begin{table}[H]
\centering
\caption{Technology versions used in the project}
\begin{tabular}{@{}ll@{}}
\toprule
\textbf{Technology} & \textbf{Version} \\
\midrule
Java                & 17 \\
Spring Boot         & 3.5.9 \\
Gradle              & 8.5 \\
Kafka               & 4.1.1 (KRaft) \\
Redis               & 7-alpine \\
MySQL               & 8.0 \\
Prometheus          & latest \\
Grafana             & latest \\
Toxiproxy           & latest \\
React               & 18 \\
Vite                & 6.x \\
Python              & 3.11 \\
Bucket4j            & 8.10.1 \\
Nginx               & alpine \\
\bottomrule
\end{tabular}
\end{table}

\section{Port Reference}
\label{app:ports}

\begin{table}[H]
\centering
\caption{Service port mapping}
\begin{tabular}{@{}lrrp{5cm}@{}}
\toprule
\textbf{Service}     & \textbf{Internal} & \textbf{External} & \textbf{Purpose} \\
\midrule
MySQL (direct)       & 3306 & —    & Internal only \\
MySQL (proxy)        & —    & 3307 & Toxiproxy testing \\
Redis Master (direct)& 6379 & 6379 & Cache R/W \\
Redis Master (proxy) & —    & 26379& Toxiproxy testing \\
Redis Replica        & 6379 & 6380 & Read-only standby \\
Kafka (internal)     & 9092 & 9092 & Event bus \\
Kafka (proxy)        & —    & 29093& Toxiproxy testing \\
Spring Boot API      & 8080 & 8080 & REST + SSE \\
React Frontend       &   80 & 3000 & Web UI \\
Grafana              & 3000 & 3001 & Dashboards \\
Prometheus           & 9090 & 9090 & Metrics \\
Toxiproxy API        & 8474 & 8666 & Toxic management \\
\bottomrule
\end{tabular}
\end{table}

\end{document}
